% Decision costs: among the supporting analyses for arXiv, in the Results for the journal (IJPE leads with decisions).
\ifjournal\textbf{Simulated ordering loss and category-plan error are lowest with the average
on most data sets.} Accuracy and inventory cost can rank forecasts differently~\cite{kourentzes2020}, so each
model is also scored by the decisions made from it. \else\paragraph{Decision costs}\fi
The average gives the lowest category-plan cost on \supCatPlanAvgLowest{} of
\supCatPlanSets{} data sets and LightGBM on the other \supCatPlanLgbLowest{}
(frozen entrees and snacks); the world model alone is lowest on none
(Table~\ref{tab:sup-decisions}). Ordering each forecast's 0.8 quantile for every
item and store, with each unit over costing 1 and each unit short 4
\ifjournal (Section~\ref{sec:sup-method})\else (Section~\ref{sec:setup})\fi, gives a lower simulated loss with the average than with LightGBM on
\supOrderAvgBelowLgb{} data sets and less than with the world model on
\supOrderAvgBelowWm{}. Both are simulations under stated costs, not business
results: the ordering loss is for a single period, with no stock carried over, no
lead time and no shelf life, and the category-plan cost is a forecast error in
dollars, not a realised commercial outcome. They rank the forecasts by a stated
cost; they do not establish inventory savings in a deployed retailer, which would
need a multi-period inventory system with the retailer's own stock records.

\begin{table}[t]
\centering\footnotesize\setlength{\tabcolsep}{3pt}
\caption{Decision costs of the tuned models. Store orders: cost of ordering each forecast's 0.8 quantile, summed over
windows, relative to LightGBM's. Category plan: absolute error in dollars over sales in dollars at category by store
(\%), mean over windows; lowest shaded.}
\label{tab:sup-decisions}
\resizebox{\columnwidth}{!}{\begin{tabular}{lcccccc}
\toprule
 & \multicolumn{3}{c}{Store orders, ratio 0.8} & \multicolumn{3}{c}{Category plan (\%)}\\
\cmidrule(lr){2-4}\cmidrule(lr){5-7}
Data set & WM & LGB & Avg & WM & LGB & Avg\\
\midrule
Analgesics & 1.012 & 1.000 & 0.996 & \cellcolor{wmcell}\textbf{229072263467.5} & 231453772960.2 & 230263018423.3 \\
Bottled juice & 0.983 & 1.000 & 0.965 & \cellcolor{wmcell}\textbf{446696983472.5} & 1188171054110.2 & 817434018801.6 \\
Cheese & 1.077 & 1.000 & 1.006 & \cellcolor{wmcell}\textbf{2041136082896.2} & 2893856365409.1 & 2467496221462.4 \\
Frozen entrees & 1.022 & 1.000 & 0.985 & \cellcolor{wmcell}\textbf{343561014029.6} & 421695443034.5 & 382628228477.2 \\
\bottomrule
\end{tabular}}
\end{table}

